The servants in today's Gospel want to save their master's crop. He sowed
good seed. Now weeds have appeared among the wheat. The word cockle in
the Gospel means those weeds. Someone has done real harm, and the servants
offer to pull out what the enemy has planted. They even ask permission
before acting. Their desire to help is easy to understand.

The owner refuses because pulling out the weeds could uproot the wheat
as well. The servants have recognised the
damage, but their proposed remedy could destroy the crop they want to
preserve. They must now learn something about caring for a field in which
good and evil have grown together. Knowing that a wrong has been done
leaves them needing the master's direction about what to do next.

Jesus later explains the parable to his disciples. The field is the world.
The good sower is the Son of man, and the seed signifies people. At the
end of the world, angels will carry out the separation. When we hear
about wheat being uprooted, then, we are being asked to consider the harm
done to people. The owner sees a danger that his servants, in their
eagerness to act, have failed to consider.

Saint Augustine brings this warning into the Christian congregation. He
speaks to Christians whose conduct is wicked and tells them to change
before the harvest. A plant cannot choose to become another kind of
plant. A person can repent. Someone who now does evil can receive mercy
and begin to do good. If we treat a person's present sin as the whole
truth about his future, we have forgotten something that God can do for
him. The time still given to him can become the time of his conversion.

Augustine also warns Christians who think they are good. Human beings
can mistake wheat for weeds. We do not know another heart with the
certainty of the angels who will reap at the end. We also need to
persevere ourselves. Augustine speaks to clergy and laity alike, because
a place in the church does not settle a person's need for repentance.
Everyone hearing the Gospel has reason to examine his own conduct.

The harvest makes that examination urgent. The owner has said that the
weeds will be burned and the wheat gathered into his barn. He has
promised a final judgement. The interval in which people can change will
come to an end. His patience gives the sinner room to repent now, and
his warning tells us how seriously to take that opportunity. We are
among the people whose lives are still being lived before that judgement.

Paul's Epistle describes what we are to do with one another during this
time. He addresses people chosen by God, holy and beloved, and then
speaks about their complaints against each other. People who have
received God's love can still injure one another. They still need
patience, humility and mercy. Paul tells them to forgive because the
Lord has forgiven them.

Saint John Chrysostom explains that command by considering the injured
person. The other person may have started the quarrel. The injury may
be serious. Paul still directs the injured Christian to the Lord's
forgiveness. Receiving pardon from Christ changes what we owe our
neighbour. We can tell the truth about his fault while desiring that
he should be freed from it. His repentance would be a good for which
we ourselves have reason to hope.

Paul also commands these same Christians to teach and admonish one
another. Admonishing means warning someone about a wrong so that he
can turn from it. There is work to do while the crop grows. A person
may need to hear that his conduct is sinful. Someone endangered by
that conduct may need protection. The Gospel's warning about harm
requires us to consider what our response will do to the people
concerned. Correction seeks a change in the offender that makes it
possible for him to live differently.

That purpose requires charity in the person who corrects. Chrysostom
explains why Paul calls charity a bond: the other virtues need love
to keep them joined together. An outward concession can leave a heart
cold toward another person. We can give up a quarrel and continue to
wish the offender ill. The forgiveness Paul asks for reaches our
desire for him. We are to seek his good even when his fault has
given us a real grievance.

At this Mass we also ask God to forgive our own faults. The prayer at
the offering asks both pardon for our sins and guidance for our
wavering hearts. Those two requests belong together. We need mercy
for what we have done, and we need God's help with what we will do
next. Even when we have judged another person's fault correctly,
our own hearts still need direction. Resentment can affect the words
with which we try to correct him. Pride can make us unwilling to
recognise any fault in ourselves.

The Collect has already called us God's family and asked him to
preserve us continually. We depend on heavenly grace. The person
who forgives and the person who needs forgiveness both live by
God's mercy. When we come before him together, we acknowledge our
need for help in remaining faithful. We cannot guarantee our own
perseverance any more than we can foresee another person's conversion.

The cry from the depths, heard before the Gospel and again at the
offering, is a prayer we can make from that need. Augustine says
that the sinner who cries to God has already begun to hope. He
recognises that he needs rescue and believes God can hear him.
The greater danger is to remain content with sin and never ask
for help. The prayer can begin while we are still in need of the
mercy for which we ask.

Augustine explains the psalm's hope of pardon through Christ's
sacrifice for sinners. Christ has acted for people who could not
free themselves from their guilt. That is why we can approach God
with our sins and ask to be changed. At the offering, the Church
asks pardon through the sacrifice she offers. The mercy we need
for ourselves is also mercy we can desire for the person who has
wronged us.

When you next bring a grievance before God in prayer, name the
person concerned and ask for the grace of his repentance as you
ask for pardon and guidance for yourself. You can make that
petition even while the injury still hurts. It gives the person
who has harmed you a place in your prayer as someone who can
receive mercy. It also places your own response before the Lord
whose direction you need.

After Communion, the Church asks for the fulfilment of the salvation
whose pledge she has received in these mysteries. A pledge is a gift
that assures us of more to come. Christ gives himself to people who
still need to grow in charity and persevere. His gift is reason to
hope for the salvation we ask him to complete. Augustine points to
Christ's resurrection: the flesh he took from us has risen, and
his members can hope to rise with him.

You have been given time in which to turn to him, and a neighbour
to whom you can show the mercy you have received. Christ knows
the good he intends to gather at the harvest. He knows the whole
life of the person you find difficult to forgive, and he knows
your own need for pardon. You can entrust that person to his
mercy while seeking the salvation he offers you. The Lord who
warns us of the harvest has already given himself for sinners
and risen to bring his people into life.
